\scanpage{6}
\mspage{4}
dans \(\mathbb P^3_k\), alors on a
\[
\mathbb Z\simeq\operatorname{Pic}(\mathbb P^3)\xrightarrow{\sim}\operatorname{Pic}(X_0),
\]
\(k\) étant un corps de car. quelconque.

2. \uline{Rappels sur la structure de} \(\pi\).

On suppose \(S\) hensélien, i.e. \(\deleted{\unsure{R|A}}V\) hensélien.

2.1. Il y a une sous-extension \(\widetilde K_{\deleted{\unsure{e|nr}}}\) de \(\overline K/K\), qui est la plus grande ss-extension non ramifiée (i.e. étale) \deleted{\unsure{dans}} sur \(V\). Si \deleted{\(\unsure{\overline K_e}\) est} \(\overline V\) est la clôture normale de \(V\) ds \(\overline K\), \(\widetilde V=\deleted{\unsure{\overline V_e}=}\overline V\cap\widetilde K_{\deleted{e}}\) la clôture normale dans \(\widetilde K_{\deleted{e}}\), alors \(\widetilde V_{\deleted{\unsure{e}}}\) est un anneau de valuation discrète tel que
\[
\begin{aligned}
&\deleted{\unsure{k|M|\mathfrak m}(\widetilde V)=\mathfrak m\widetilde V,}\\
&\widetilde V\otimes_V k=\bar k
\end{aligned}
\]
est une \deleted{extension \unsure{al}} clôture séparable de \(k\) ; En d'autres termes, \(\widetilde V\) est une \uline{hensélisation stricte} de \(V\) p.r. à \deleted{\unsure{une}} la \deleted{\unsure{co}} clôture séparable \(\bar k\) de \(k\), et \(\widetilde K_{\deleted{e}}\) est le corps des fractions de \(\widetilde V\). On posera \inserted{\(\left\{\begin{aligned}\widetilde S&=\operatorname{Spec}\widetilde V=\text{hensélisé strict de }S\\ \bar s&=\operatorname{Spec}\bar k.\end{aligned}\right.\)}

\margin{\(\circ\mathrel{-}\)} 2.3. On se rappellera que \deleted{\unsure{de sorte|de même}} la donnée de la clôture séparable \(\overline K\) de \(K\) définit canoniquement une clôture séparable \(\bar k\) de \(k\). \deleted{\unsure{Cett|Cette}} Si \(S'\) est un

\deleted{Pour tout \deleted{\unsure{entier le}} entier \(n\) \deleted{premier à} inversible dans \(k\), on a
\[
{}_{n}\bar k^{*}\xleftarrow[\deleted{\unsure{p}}]{\sim}{}_{n}\widetilde V^{*}\xrightarrow{\sim}\overline K^{*}\tag*{2.1.1.}
\]
d'où un isomorphisme canonique
\[
\mu_n(\overline K)\simeq\mu_n(\bar k).
\]
Ces isom sont compatibles avec les morphismes de}
% archive sheet number: 5
